\subsection{Additional source qualifications}\label{app:additionalevidence}

The details below qualify specific uses of the literature in
Section~\ref{sec:background}; they are not additional results from this study.

\begin{table}[!htbp]
\centering\footnotesize
\caption{Source-specific qualifications relevant to the proposed measurements.}
\label{tab:sourcequalifications}
\begin{tabularx}{\textwidth}{@{}L{2.2cm}YY@{}}
\toprule
Source & Additional detail & Consequence for interpretation \\
\midrule
Sinha et al.~\cite[Sec.~3.2, Figs.~5--6]{sinha2025illusion} &
The artificial-history experiment varies erroneous prior updates.
Qwen3 with thinking enabled avoids the reported effect in that setup,
whereas scaling non-thinking models does not remove it. &
Earlier-error sensitivity depends on configuration; it is not a universal
property of all models or a measured repository-decay rate. \\
Pan et al.~\cite[Sec.~5.1, Fig.~4]{pan2026swemaintain} &
In the \texttt{parse\_crontab} episode, a local Python 3 change satisfies
current tests but omits module-wide compatibility aliases later code uses. &
An initially successful patch can leave a latent dependency failure.
This complements, rather than repeats, the line-terminator example in the main text. \\
Gloaguen et al.~\cite[Secs.~4.2--4.4]{gloaguen2026agentsmd} &
File length does not explain success in their analysis, and generated
context helps when other documentation is removed. &
Instruction length alone is not an adequate exposure or burden measure;
available alternative documentation matters. \\
Jiang et al.~\cite[Secs.~5.2--5.5]{jiang2026skills} &
Skill misapplication and missing runtime verification remain failure modes.
The retrieval diagnostic is separate from downstream execution. &
Lower exact-match access alongside sustained task success does not isolate
an effect of retrieval; accessed, applied and sufficient guidance differ. \\
\bottomrule
\end{tabularx}
\end{table}

The provenance of merge-outcome studies also matters. Peralta et al.
start with 9,799 human-reviewed agent PRs and manually inspect 717.
Their 35.7\% clear-failure share among rejections and 15.4\% explicit
reviewer-involvement share among merges describe that manual sample
~\cite{peralta2026merged}. Murphy-Hill et al.'s roughly 24\% higher merge
count among adopters is an observational rollout estimate, not delivered
value~\cite{murphyhill2026agents}. AIDev's 932,791 agent-authored PRs from
five agents across 116,211 repositories provide a population for research,
not a count of accepted user outcomes~\cite{li2026aidev}.

Pavlenko et al.'s World Model offers an additional context-management
design: preserve raw evidence, derive rebuildable views and retrieve compact,
task-specific context with provenance and access controls
~\cite{pavlenko2026factory}. That architecture is an alternative to relying
on an expanding read path; its effect on the AMA case is untested.
